\documentclass[11pt,a4paper]{article}
\usepackage[margin=25mm]{geometry}
\usepackage[T1]{fontenc}
\usepackage{hyperref}
\newcommand{\code}[1]{\texttt{\detokenize{#1}}}
\title{Blame: Refactoring Summary}
\author{Blame Engineering}
\date{October 7, 2026}
\begin{document}
\maketitle

\section{Purpose}
The first refactoring pass separates responsibilities in the largest agent,
workspace, backlog, editor, conversation, and preview files. Public entry points
and existing product behavior are retained. The work creates smaller boundaries
for further maintenance and open-source preparation; it does not change the
licensing model or implement subscription-independent self-hosting.

\section{Extracted responsibilities}
\subsection{Agent SDK adapter}
\code{cline-runtime.ts} remains the session coordinator.
\code{cline-runtime-events.ts} owns SDK event normalization, terminal event
interpretation, duplicate signatures, transport failure classification, and
usage normalization. \code{cline-runtime-prompts.ts} builds development and
generation prompts and attachment context. \code{cline-runtime-configuration.ts}
owns backend options, tool policies, and SDK configuration builders.
\code{cline-runtime-tools.ts} assembles harness, preview, completion, and
specialist tools through the existing ports and maps approval requests.
\code{cline-runtime-log.ts} owns log paths, serialization, and file appends;
the coordinator still owns write sequencing.
\code{cline-initialization.ts} owns structured initialization and code-practice
submission, parsing, and correction. Existing exported helpers remain available
through the original runtime module.

\subsection{Desktop workspace operations}
\code{workspaces/workspace-git-service.ts} owns Git status, log, diff, conflict,
staging, commit, push, and repository fingerprint operations.
\code{workspace-ipc.ts} retains transport handlers, write barriers, and workspace
orchestration. Lower-level command and runtime adapters remain in the existing
Git and workspace runtime modules. The extracted service does not register IPC
handlers or import the IPC coordinator.

\subsection{React composition}
\code{markdown-preview.tsx} owns Markdown preview blocks, streaming reveals,
tables, inline formatting, and file links. Editing, attachments, selection, and
input controls remain in \code{markdown-editor.tsx}.
\code{native-preview.tsx} owns native browser panels, lifecycle hooks, retained
sessions, bounds synchronization, screenshots, and ambilight. The story preview
component composes the runner, native browser, iframe, controls, and logs.
\code{agent-activity-presentation.ts} interprets tool detail payloads and maps
edited or inspected files into display data. Conversation components retain
timeline and interaction responsibilities.

\subsection{Backlog and conversation persistence}
\code{user-story-agent-updates.ts} owns activity coalescing, metadata filtering,
and event-to-story/run updates. Authentication error notifications stay in the
page instead of the data mapper. \code{backlog-title-generation.ts} owns queued
title generation and run-configuration selection.
Both backlog and chat use \code{debounced-update-queue.ts}: independent keys,
replacement of pending updates, cancellation before terminal persistence, and
cleanup on unmount. A stale timer cannot consume a newer entry. The backlog
terminal outbox and chat snapshot finalization paths retain their existing
storage and retry semantics.

\section{Behavioral checks}
Queue regression tests cover coalescing independent targets, cancelled terminal
writes, stale callbacks, unmount cleanup, and reentrant persistence. Story
mapping tests cover activity ordering and replacement, authentication pause,
and successful review transitions. Existing application, agent-runtime, and
workspace tests are used to check the extracted boundaries.

The validation snapshot includes 258 passing application tests, 118 passing
focused runtime tests, and 45 passing focused workspace tests. The runtime
TypeScript build and frontend bundle succeed. App and desktop type checking
still report the same shared runtime errors as the pre-change baseline. Broader
suites expose command timing, formatter, and archive synchronization failures;
baseline probes reproduce failures in those paths. The six targeted files now
meet their existing line ceilings, while other repository files still exceed
their budgets. Linux tests do not
establish physical macOS/Lima, Windows/WSL, Electron browser, or live cloud
provider coverage. The repository-wide line-budget check remains a separate
constraint, including legacy files outside this pass. The native LaTeX
editor/compiler requests did not return a result in this environment: sources
have been checked for balanced braces and document environments, but PDF
compilation remains unverified.

\section{Remaining work}
The runtime session coordinator and workspace coordinator are still large.
Future passes should isolate per-session lifecycle state and teardown, then
workspace setup and recovery workflows. Markdown DOM selection and editor
interaction, conversation rendering, and backlog orchestration remain further
extraction candidates. Each pass should preserve cancellation, write ownership,
terminal finalization, preview retention, and acknowledgment ordering.

The documentation directory now contains a small set of English, standalone
LaTeX summaries. Historical specifications and benchmark outputs are removed
from the current tree and remain recoverable from Git history. Development
fixtures and production resources outside that directory are retained.
\end{document}
